\begin{theorem}[ResidueOfTendsTo]
\label{ResidueOfTendsTo}
\lean{ResidueOfTendsTo}
% at /home/runner/work/trough/trough/library/PrimeNumberTheoremAnd/ZetaBounds.lean:41.0-159.33
\leanok
If a function $f$ is holomorphic in a neighborhood of $p$ and
  $\lim_{s\to p} (s-p)f(s) = A$, then
  $f(s) = \frac{A}{s-p} + O(1)$ near $p$.
\end{theorem}
\begin{proof}
\uses{existsDifferentiableOn_of_bddAbove}
\leanok
The function $(s - p)\cdot f(s)$ bounded, so by Theorem
  \ref{existsDifferentiableOn_of_bddAbove}, there is a holomorphic function, $g$, say, so that
  $(s-p)f(s) = g(s)$ in a neighborhood of $s=p$, and $g(p)=A$. Now because $g$ is holomorphic,
  near $s=p$, we have $g(s)=A+O(s-p)$. Then when you divide by $(s-p)$, you get
  $f(s) = A/(s-p) + O(1)$.
\end{proof}
